% Lab Tracker — Специфікація вимог до програмного забезпечення (SRS)
% Збирається pdflatex (або в Overleaf). Одне речення — один рядок, щоб diff був читабельним.
\documentclass[12pt,a4paper]{article}
\usepackage[utf8]{inputenc}
\usepackage[T2A]{fontenc}
\usepackage[ukrainian]{babel}
\usepackage[margin=2cm]{geometry}
\usepackage{longtable,booktabs,array}
\usepackage{enumitem}
\usepackage[hidelinks]{hyperref}

\newcommand{\code}[1]{\texttt{#1}}
\newcommand{\term}[1]{\textit{#1}}
% Вимога: \req{ID}{Назва}
\newcommand{\req}[2]{\subsubsection*{#1. #2}\label{#1}}
\newcommand{\ac}[1]{\par\smallskip\noindent\textbf{Критерій прийняття (AC-#1).}~}
\newcommand{\bdd}[1]{\par\noindent\textbf{BDD:}~\code{#1}}
\newcommand{\src}[1]{\par\noindent\textit{Джерело:}~#1}

\title{Lab Tracker\\\large Специфікація вимог до програмного забезпечення (SRS)}
\author{Версія 1.0 (базова)}
\date{2026-11-19}

\begin{document}
\maketitle

\begin{center}
\begin{tabular}{ll}
Статус & Базова специфікація (baseline), затверджується на SPEC-GATE \\
Вхідні артефакти & \code{spec/vision.md}, \code{logs/agent.md}, \code{logs/conflict-statuses.md} \\
Протокол узгодження & \code{logs/sked-dialog.md} \\
BDD-сценарії & \code{spec/features/*.feature} \\
Простежуваність & \code{spec/traceability.md} \\
\end{tabular}
\end{center}

\tableofcontents

% =====================================================================
\section{Вступ}

\subsection{Призначення документа}
Документ визначає вимоги до версії 1.0 системи Lab Tracker.
Він є еталоном для декомпозиції на задачі, створення тестів і реалізації в подальшому SDD-циклі.
Поведінка, яку не описано в цьому документі, не повинна реалізовуватися без зміни специфікації (адендуму).
Ключові слова \textbf{повинна} та \textbf{не повинна} позначають обов'язкові вимоги.

\subsection{Призначення системи}
Lab Tracker --- консольний застосунок, у якому студент веде облік своїх лабораторних робіт з різних дисциплін.
Система розв'язує проблему розпорошеності інформації про терміни та стан робіт між Moodle, месенджерами та нотатками, через яку студент пропускає дедлайни й забуває, які роботи ще не захищено.

\subsection{Межі системи}

\paragraph{Входить у версію 1.0.}
\begin{itemize}[nosep]
  \item Облік дисциплін і лабораторних робіт (додавання, перегляд, редагування, видалення).
  \item Відстеження стану роботи за життєвим циклом (розділ~\ref{sec:lifecycle}), зокрема пропуск етапів і скасування останнього переходу.
  \item Облік балів: максимальний бал роботи та отримані бали при захисті.
  \item Обчислювана ознака прострочення.
  \item Перегляд найближчих дедлайнів і підсумок дисципліни.
  \item Щоденне нагадування при першій команді дня.
  \item Локальне зберігання даних у файлі.
\end{itemize}

\paragraph{Не входить у версію 1.0.}
\begin{itemize}[nosep]
  \item Роль викладача та будь-які інші ролі, крім студента; кілька користувачів; автентифікація.
  \item Інтеграція з Moodle та іншими зовнішніми сервісами; будь-який мережевий обмін.
  \item Графічний, вебінтерфейс або мобільний застосунок.
  \item Синхронізація між пристроями.
  \item Сповіщення поза запуском команди (push, email, планувальник ОС).
  \item Експорт у CSV або інші формати.
  \item Архівування дисциплін і поділ на семестри.
  \item Історія переходів станів глибше за один крок; кошик і відновлення видалених даних.
\end{itemize}

\subsection{Глосарій}\label{sec:glossary}
Терміни глосарію та їхні коди є єдиними назвами для специфікації, тікетів, тестів, коду й документації.

\begin{longtable}{>{\raggedright}p{0.22\textwidth} >{\raggedright}p{0.15\textwidth} p{0.55\textwidth}}
\toprule
\textbf{Термін} & \textbf{Код} & \textbf{Визначення} \\
\midrule
\endhead
Студент & --- & Єдиний користувач системи; веде облік власних робіт. \\
Дисципліна & \code{discipline} & Навчальна дисципліна, до якої належать роботи. Ідентифікується назвою. \\
Лабораторна робота (робота) & \code{lab} & Завдання в межах дисципліни, стан якого відстежує студент. Належить рівно одній дисципліні. \\
Номер роботи & \code{number} & Ціле число $\ge 1$, унікальне в межах дисципліни. Разом із назвою дисципліни однозначно визначає роботу. \\
Дедлайн & \code{deadline} & Календарна дата (без часу), до якої \textbf{включно} роботу треба здати. Необов'язковий атрибут. \\
Сьогодні & \code{today} & Поточна календарна дата за системним годинником і часовим поясом комп'ютера користувача. \\
Стан роботи & \code{status} & Одне зі значень: \code{not\_started}, \code{in\_progress}, \code{submitted}, \code{returned}, \code{defended} (розділ~\ref{sec:lifecycle}). \\
Не розпочата & \code{not\_started} & Роботу додано, виконання не почалося. Початковий стан. \\
Виконується & \code{in\_progress} & Студент працює над роботою. \\
Здана & \code{submitted} & Звіт надіслано викладачу, рішення про захист ще немає. \textbf{Увага:} у розмовній мові «здав лабу» часто означає «захистив»; у системі «здана» означає лише передачу звіту. \\
Повернута на доопрацювання (повернута) & \code{returned} & Викладач не прийняв здану роботу; її треба виправити й здати повторно. \\
Захищена & \code{defended} & Роботу прийнято, виставлено бали. Синонім «оцінена» в системі \textbf{не використовується}. \\
Перехід & --- & Зміна стану роботи командою студента. \\
Пропуск етапів & --- & Перехід вперед через один чи кілька проміжних станів (наприклад, \code{not\_started} $\to$ \code{defended}). \\
Попередній стан & \code{previous\_status} & Стан роботи безпосередньо перед останнім переходом; використовується для скасування. \\
Скасування переходу & \code{undo} & Повернення роботи в попередній стан. Глибина --- один крок. \\
Максимальний бал & \code{max\_score} & Ціле число $\ge 1$ --- найбільша кількість балів за роботу. \\
Отримані бали & \code{score} & Ціле число від 0 до максимального балу; існує лише в стані \code{defended}. \\
Прострочена робота & \code{overdue} & \textbf{Обчислювана ознака, а не стан}: дедлайн задано, дедлайн $<$ сьогодні, стан \code{not\_started} або \code{in\_progress}. \\
Вікно найближчих дедлайнів & --- & Дати від сьогодні до сьогодні $+ 7$ днів включно. \\
Підсумок дисципліни & --- & Кількість захищених робіт від загальної кількості та сума балів захищених робіт. \\
Щоденне нагадування & --- & Короткий рядок про прострочені й термінові роботи, що виводиться при першій команді календарного дня. \\
Файл даних & --- & Локальний файл, у якому зберігаються всі дані системи. \\
\bottomrule
\end{longtable}

\subsection{Посилання}
\begin{itemize}[nosep]
  \item \code{spec/vision.md} --- початковий задум системи.
  \item \code{logs/sked-dialog.md} --- протокол SKED-діалогу та рішення людини.
  \item IEEE Std 830-1998 --- рекомендована практика специфікації вимог (адаптовано).
\end{itemize}

% =====================================================================
\section{Загальний опис}

\subsection{Перспектива продукту}
Lab Tracker --- автономний застосунок командного рядка без зовнішніх залежностей від мережевих сервісів.
Усі дані зберігаються у файлі даних на комп'ютері користувача.

\subsection{Узагальнені функції}
Облік дисциплін і робіт; керування станом роботи; облік балів; перегляд прострочених робіт і найближчих дедлайнів; підсумок дисципліни; щоденне нагадування.

\subsection{Користувачі}
Студент, який упевнено користується командним рядком.
Інших ролей у версії 1.0 немає.

\subsection{Середовище функціонування}
Windows 10/11 та Linux; інтерпретатор Python версії 3.11 або новішої.

\subsection{Обмеження проєкту}
\begin{itemize}[nosep]
  \item C1. Інтерфейс --- лише командний рядок (команди з аргументами).
  \item C2. Мова реалізації --- Python $\ge$ 3.11.
  \item C3. Дані зберігаються локально у файлі формату JSON; база даних і мережевий обмін не використовуються.
  \item C4. Один користувач; автентифікації немає.
\end{itemize}

\subsection{Припущення та залежності}
\begin{itemize}[nosep]
  \item A1. Робота належить рівно одній дисципліні; кількість робіт у дисципліні не обмежена.
  \item A2. Нотатки --- довільний текст; посилання на завдання не перевіряється на доступність.
  \item A3. Системний годинник комп'ютера налаштовано правильно.
  \item A4. Видалення роботи остаточне.
  \item A5. Очікуваний обсяг даних --- до 20 дисциплін і 300 робіт.
\end{itemize}

% =====================================================================
\section{Модель даних}

\paragraph{Дисципліна.}
Атрибут: назва --- непорожній рядок до 100 символів, унікальний без урахування регістру.

\paragraph{Лабораторна робота.}
\begin{longtable}{>{\raggedright}p{0.2\textwidth} p{0.14\textwidth} p{0.56\textwidth}}
\toprule
\textbf{Атрибут} & \textbf{Обов'язковий} & \textbf{Правила} \\
\midrule
Дисципліна & так & Існуюча дисципліна. \\
Номер & так & Ціле $\ge 1$, унікальне в межах дисципліни. \\
Назва & так & Непорожній рядок до 200 символів. \\
Дедлайн & ні & Дата; може бути змінена або видалена. \\
Посилання на завдання & ні & Довільний рядок. \\
Нотатки & ні & Довільний текст. \\
Максимальний бал & ні* & Ціле $\ge 1$. *Обов'язковий для переходу в \code{defended}. \\
Отримані бали & ні* & Ціле від 0 до максимального балу. *Обов'язкові в стані \code{defended}, відсутні в інших станах. \\
Стан & так & Значення зі списку станів; початково \code{not\_started}. \\
Попередній стан & ні & Встановлюється при кожному переході; очищується після скасування. \\
\bottomrule
\end{longtable}

% =====================================================================
\section{Життєвий цикл лабораторної роботи}\label{sec:lifecycle}

Робота перебуває рівно в одному стані.
Основна послідовність: \code{not\_started} $\to$ \code{in\_progress} $\to$ \code{submitted} $\to$ \code{defended}.
Здану роботу може бути повернуто на доопрацювання (\code{returned}), після чого її здають повторно.
Перехід вперед може пропускати проміжні стани.
Стан \code{returned} досяжний лише зі стану \code{submitted}.
Стан \code{defended} є кінцевим: вийти з нього можна лише скасуванням переходу.

\begin{center}
\begin{tabular}{lll}
\toprule
\textbf{Поточний стан} & \textbf{Дозволені переходи} & \textbf{Примітка} \\
\midrule
\code{not\_started} & \code{in\_progress}, \code{submitted}, \code{defended} & пропуск дозволено \\
\code{in\_progress} & \code{submitted}, \code{defended} & пропуск дозволено \\
\code{submitted} & \code{returned}, \code{defended} & \\
\code{returned} & \code{submitted}, \code{defended} & пропуск дозволено \\
\code{defended} & --- & лише скасування \\
\bottomrule
\end{tabular}
\end{center}

Усі переходи, не наведені в таблиці (зокрема перехід у той самий стан і будь-який перехід назад), заборонені.
Повернення назад можливе лише скасуванням останнього переходу (REQ-F-012).

% =====================================================================
\section{Функціональні вимоги}

\subsection{Дисципліни}

\req{REQ-F-001}{Додавання дисципліни}
Система повинна дозволяти додати дисципліну із заданою назвою.
Якщо дисципліна з такою назвою (без урахування регістру) вже існує, система повинна відхилити команду.
\ac{F-001} Після додавання дисципліна присутня в списку дисциплін; спроба додати дисципліну «ООП», коли існує «ооп», завершується повідомленням про помилку, а кількість дисциплін не змінюється.
\bdd{SC-01, SC-02}

\req{REQ-F-002}{Перегляд списку дисциплін}
Система повинна виводити список дисциплін, упорядкований за назвою, із кількістю робіт у кожній.
\ac{F-002} Для дисциплін «Бази даних» (2 роботи) та «Алгоритми» (0 робіт) виводиться спочатку «Алгоритми --- 0», потім «Бази даних --- 2».

\req{REQ-F-003}{Видалення дисципліни}
Система повинна дозволяти видалити дисципліну.
Якщо дисципліна містить роботи, система повинна запитати підтвердження і в разі згоди видалити дисципліну разом з усіма її роботами; без підтвердження дані не змінюються.
\ac{F-003} Дисципліна без робіт видаляється без запиту; дисципліна з роботами видаляється лише після підтвердження або з прапорцем \code{--yes}; у разі відмови дисципліна і роботи залишаються.
\bdd{SC-03}

\subsection{Лабораторні роботи}

\req{REQ-F-004}{Додавання роботи}
Система повинна дозволяти додати роботу до існуючої дисципліни із зазначенням номера та назви, а також, за бажанням, дедлайну, посилання, нотаток і максимального балу.
Нова робота повинна отримувати стан \code{not\_started}.
Система повинна відхиляти роботу з номером, який уже існує в цій дисципліні.
\ac{F-004} Після додавання робота відображається в списку робіт зі станом «Не розпочата»; повторне додавання роботи №1 до тієї ж дисципліни відхиляється; робота №1 в іншій дисципліні додається успішно.
\bdd{SC-04, SC-05}

\req{REQ-F-005}{Редагування роботи}
Система повинна дозволяти змінити назву, номер, дедлайн (зокрема видалити його), посилання, нотатки та максимальний бал роботи з дотриманням правил моделі даних.
Для захищеної роботи новий максимальний бал не може бути меншим за отримані бали.
\ac{F-005} Після зміни дедлайну картка роботи показує нову дату; спроба встановити максимальний бал 5 для захищеної роботи з 8 балами відхиляється.
\bdd{SC-06}

\req{REQ-F-006}{Видалення роботи}
Система повинна дозволяти остаточно видалити роботу.
\ac{F-006} Після видалення робота відсутня у списку робіт і підсумку дисципліни.

\req{REQ-F-007}{Перегляд списку робіт}
Система повинна виводити список робіт (усіх або однієї дисципліни), упорядкований за назвою дисципліни, а потім за номером, із номером, назвою, дедлайном, станом і позначкою прострочення.
\ac{F-007} Для прострочених робіт у рядку присутня позначка «ПРОСТРОЧЕНО»; фільтр за дисципліною виводить лише її роботи.

\req{REQ-F-008}{Перегляд картки роботи}
Система повинна виводити всі атрибути роботи, зокрема стан, позначку прострочення, максимальний і отримані бали.
\ac{F-008} Картка захищеної роботи містить рядок виду «Бали: 8 з 10».

\subsection{Стан роботи}

\req{REQ-F-009}{Зміна стану роботи}
Система повинна дозволяти переводити роботу в інший стан відповідно до таблиці переходів (розділ~\ref{sec:lifecycle}) і відхиляти недозволені переходи.
При кожному успішному переході система повинна запам'ятовувати попередній стан.
\ac{F-009} Переходи \code{not\_started} $\to$ \code{submitted} та \code{in\_progress} $\to$ \code{defended} (з балами) виконуються; переходи \code{submitted} $\to$ \code{in\_progress} та \code{defended} $\to$ \code{submitted} відхиляються з повідомленням про недозволений перехід.
\bdd{SC-07, SC-08}
\src{\code{spec/vision.md} §5; рішення SKED 2.3}

\req{REQ-F-010}{Обмеження переходу в стан «Повернута»}
Система повинна дозволяти перехід у стан \code{returned} лише зі стану \code{submitted}.
\ac{F-010} Перехід \code{in\_progress} $\to$ \code{returned} відхиляється; перехід \code{submitted} $\to$ \code{returned} виконується.
\bdd{SC-09}
\src{рішення SKED 3.1}

\req{REQ-F-011}{Захист роботи з балами}
Перехід у стан \code{defended} повинен вимагати введення отриманих балів --- цілого числа від 0 до максимального балу роботи.
Якщо максимальний бал роботи не задано, система повинна відхилити перехід і вказати, що спочатку треба задати максимальний бал.
\ac{F-011} Захист із балами 8 при максимумі 10 виконується; захист без балів, з балами 11 при максимумі 10, з балами $-1$ або для роботи без максимального балу відхиляється, і стан роботи не змінюється.
\bdd{SC-10, SC-11, SC-12}
\src{\code{spec/vision.md} §7 Q2; рішення SKED 2, 2.1}

\req{REQ-F-012}{Скасування останнього переходу}
Система повинна дозволяти скасувати останній перехід роботи, повертаючи її в попередній стан, зокрема після пропуску етапів.
Глибина скасування --- один крок: після скасування повторне скасування неможливе до наступного переходу.
При скасуванні переходу в стан \code{defended} отримані бали повинні видалятися.
Якщо попереднього стану немає, система повинна відхилити команду.
\ac{F-012} Після переходу \code{not\_started} $\to$ \code{submitted} і скасування робота знову має стан \code{not\_started}; повторне скасування відхиляється; після скасування захисту в картці роботи немає балів.
\bdd{SC-13, SC-14}
\src{рішення SKED 6, 2.2}

\subsection{Дедлайни та огляд}

\req{REQ-F-013}{Ознака прострочення}
Система повинна вважати роботу простроченою тоді й лише тоді, коли дедлайн задано, дедлайн раніше за сьогоднішню дату, а стан роботи --- \code{not\_started} або \code{in\_progress}.
Прострочення не є станом і не обмежує переходи: прострочену роботу можна здати.
\ac{F-013} За сьогоднішньої дати 2026-11-20: робота в стані \code{in\_progress} з дедлайном 2026-11-19 прострочена; з дедлайном 2026-11-20 --- ні; робота в стані \code{returned} або \code{submitted} з дедлайном 2026-11-19 --- ні; робота без дедлайну --- ні.
\bdd{SC-15, SC-16}
\src{\code{spec/vision.md} §7 Q1; рішення SKED 1, 3.2}

\req{REQ-F-014}{Найближчі дедлайни}
Система повинна виводити: (1) окремим блоком угорі --- усі прострочені роботи; (2) роботи, дедлайн яких потрапляє у вікно найближчих дедлайнів, а стан відмінний від \code{defended}.
Обидва блоки впорядковуються за дедлайном.
\ac{F-014} За сьогоднішньої дати 2026-11-20: роботи з дедлайнами 2026-11-20 і 2026-11-27 входять у вікно, з дедлайном 2026-11-28 --- ні; захищена робота з дедлайном 2026-11-22 не виводиться; прострочена робота виводиться в блоці прострочених.
\bdd{SC-17, SC-18}
\src{\code{spec/vision.md} §4.4; рішення SKED 7}

\req{REQ-F-015}{Підсумок дисципліни}
Система повинна виводити для кожної дисципліни (або для однієї вказаної): кількість захищених робіт від загальної кількості робіт, а також суму отриманих балів і суму максимальних балів захищених робіт.
\ac{F-015} Для дисципліни з 3 роботами, з яких захищено 2 (8 з 10 і 9 з 10), виводиться «Захищено 2 з 3; бали 17 з 20».
\bdd{SC-19}
\src{\code{spec/vision.md} §4.5; рішення SKED 8, 3.3}

\req{REQ-F-016}{Щоденне нагадування}
При першій команді календарного дня система повинна перед результатом команди вивести один рядок із кількістю прострочених робіт і кількістю робіт зі станом, відмінним від \code{defended}, дедлайн яких --- сьогодні або завтра.
Якщо обидві кількості дорівнюють нулю, рядок не виводиться.
Наступні команди того самого дня нагадування не виводять.
\ac{F-016} Перша команда дня при 1 простроченій роботі та 2 роботах із дедлайном завтра виводить рядок «Нагадування: прострочено --- 1, дедлайн сьогодні/завтра --- 2»; друга команда того ж дня його не виводить; перша команда наступного дня виводить знову.
\bdd{SC-20, SC-21}
\src{рішення SKED 12, 2.4}

\subsection{Зберігання даних і помилки}

\req{REQ-F-017}{Збереження даних}
Система повинна зберігати результат кожної успішної команди, що змінює дані, у файл даних до завершення команди.
\ac{F-017} Робота, додана однією командою, відображається в списку робіт, виведеному наступною командою (окремим запуском програми).

\req{REQ-F-018}{Перший запуск}
Якщо файл даних не існує, система повинна працювати з порожнім набором даних і створити файл при першій зміні.
\ac{F-018} На чистій системі команда перегляду списку дисциплін завершується успішно з порожнім результатом.

\req{REQ-F-019}{Пошкоджений файл даних}
Якщо файл даних неможливо прочитати або він не відповідає формату, система повинна вивести повідомлення з шляхом до файлу, завершитися з кодом 2 і не змінювати файл.
\ac{F-019} Після запису у файл даних некоректного JSON будь-яка команда завершується з кодом 2, а вміст файлу побайтово не змінюється.
\bdd{SC-22}
\src{рішення SKED 10}

\req{REQ-F-020}{Обробка некоректного введення}
Система повинна відхиляти команду з некоректними або неповними аргументами, виводити причину й завершуватися з кодом 1 без зміни даних.
\ac{F-020} Додавання роботи з номером «abc» або датою «2026-13-01» завершується з кодом 1, а файл даних не змінюється.

% =====================================================================
\section{Нефункціональні вимоги}

\req{REQ-NF-001}{Продуктивність}
Будь-яка команда повинна завершуватися не довше ніж за 1 с при обсязі даних 20 дисциплін і 300 робіт на комп'ютері з SSD.
\ac{NF-001} Автоматизований тест на згенерованому наборі 20/300 вимірює час виконання кожної команди; жоден замір не перевищує 1 с.

\req{REQ-NF-002}{Цілісність даних при збоях}
Аварійне завершення програми під час збереження не повинно пошкоджувати файл даних: після збою файл містить або попередню, або нову версію даних повністю.
\ac{NF-002} Тест, що перериває збереження після часткового запису, підтверджує, що файл даних успішно читається і містить попередню версію даних.

\req{REQ-NF-003}{Переносимість}
Система повинна працювати у Windows 10/11 та Linux з Python 3.11 або новішим без зміни коду.
\ac{NF-003} Набір автоматизованих тестів проходить на Windows і на Linux з Python 3.11.

\req{REQ-NF-004}{Автономність і приватність}
Система не повинна виконувати мережевих з'єднань; дані користувача зберігаються лише у файлі даних.
\ac{NF-004} Огляд коду не виявляє використання мережевих модулів; усі тести проходять без доступу до мережі.

\req{REQ-NF-005}{Мова та формати}
Усі повідомлення для користувача повинні бути українською мовою.
Дати вводяться у форматі \code{РРРР-ММ-ДД} і виводяться у форматі \code{ДД.ММ.РРРР}.
\ac{NF-005} Тести виводу команд не містять англомовних повідомлень (крім кодів станів і назв команд); дата 2026-11-20 виводиться як 20.11.2026.

\req{REQ-NF-006}{Зручність використання}
Кожна команда повинна мати довідку (\code{--help}) з описом аргументів і прикладом.
Кожне повідомлення про помилку повинно містити причину та підказку щодо виправлення.
\ac{NF-006} Для кожної команди \code{--help} виводить опис і щонайменше один приклад; повідомлення про недозволений перехід перелічує дозволені переходи з поточного стану.

\req{REQ-NF-007}{Супровідність}
Назви станів і сутностей у коді та файлі даних повинні збігатися з кодами глосарію (розділ~\ref{sec:glossary}).
Кожен дозволений і кожен заборонений перехід із таблиці розділу~\ref{sec:lifecycle} повинен бути покритий автоматизованим тестом.
\ac{NF-007} У коді немає ідентифікаторів станів поза переліком глосарію; тести перебирають усі 25 пар (поточний стан, цільовий стан) і перевіряють результат за таблицею.

\req{REQ-NF-008}{Тестованість}
Сьогоднішню дату та шлях до файлу даних повинно бути можливо задати через змінні середовища \code{LABTRACKER\_TODAY} і \code{LABTRACKER\_DATA} без зміни коду.
\ac{NF-008} Запуск з \code{LABTRACKER\_TODAY=2026-11-20} обчислює прострочення відносно 20.11.2026; запуск з \code{LABTRACKER\_DATA} використовує вказаний файл.

% =====================================================================
\section{Вимоги до інтерфейсів}

\subsection{Інтерфейс командного рядка}
Програма запускається як \code{labtracker <команда> [аргументи]}.
Роботу ідентифікують назва дисципліни та номер роботи.

\begin{longtable}{>{\raggedright}p{0.5\textwidth} p{0.2\textwidth} p{0.2\textwidth}}
\toprule
\textbf{Команда} & \textbf{Призначення} & \textbf{Вимоги} \\
\midrule
\endhead
\code{discipline add <назва>} & додати дисципліну & REQ-F-001 \\
\code{discipline list} & список дисциплін & REQ-F-002 \\
\code{discipline remove <назва> [--yes]} & видалити дисципліну & REQ-F-003 \\
\code{lab add <дисц.> <номер> <назва> [--deadline Д] [--link L] [--notes N] [--max M]} & додати роботу & REQ-F-004 \\
\code{lab edit <дисц.> <номер> [--title T] [--number N] [--deadline Д|none] [--link L] [--notes N] [--max M]} & редагувати роботу & REQ-F-005 \\
\code{lab remove <дисц.> <номер>} & видалити роботу & REQ-F-006 \\
\code{lab list [--discipline <дисц.>]} & список робіт & REQ-F-007 \\
\code{lab show <дисц.> <номер>} & картка роботи & REQ-F-008 \\
\code{status <дисц.> <номер> <стан> [--score S]} & змінити стан & REQ-F-009--011 \\
\code{undo <дисц.> <номер>} & скасувати перехід & REQ-F-012 \\
\code{upcoming} & прострочені й найближчі & REQ-F-013, 014 \\
\code{summary [<дисц.>]} & підсумок & REQ-F-015 \\
\bottomrule
\end{longtable}

\paragraph{Коди завершення.} 0 --- успіх; 1 --- помилка введення або порушення правил; 2 --- помилка файлу даних.

\subsection{Інтерфейс файлу даних}
\begin{itemize}[nosep]
  \item Розташування за замовчуванням: \code{\textasciitilde/.labtracker/data.json}; перевизначається змінною \code{LABTRACKER\_DATA}.
  \item Формат: JSON у кодуванні UTF-8; кореневий об'єкт містить поле \code{format\_version} зі значенням 1.
  \item Файл також зберігає дату останнього показаного щоденного нагадування.
\end{itemize}

% =====================================================================
\section{Критерії прийняття версії 1.0}
Версію 1.0 приймають, якщо:
\begin{enumerate}[nosep]
  \item для кожної вимоги REQ-F і REQ-NF виконано її критерій прийняття;
  \item усі BDD-сценарії \code{spec/features/*.feature} проходять;
  \item матриця простежуваності \code{spec/traceability.md} не містить вимог без тестів.
\end{enumerate}

% =====================================================================
\section{Відкриті питання}
На момент затвердження базової специфікації відкритих питань, що блокують реалізацію, немає.
Питання Q1 і Q2 з \code{spec/vision.md} закрито вимогами REQ-F-013 та REQ-F-011.
Результати незалежної перевірки --- у \code{logs/spec-review.md}.

\end{document}
